\documentclass[11pt]{article}
\usepackage{amsmath,amssymb}
\usepackage[margin=1in]{geometry}
\usepackage{booktabs}
\usepackage{subcaption}
\usepackage{tikz}
\usepackage{pgfplots}
\pgfplotsset{compat=1.18}

\newcommand{\schematicpair}[5]{%
\begin{figure}[htbp]
\centering
\begin{subfigure}[t]{0.48\textwidth}
\centering
\begin{tikzpicture}
\begin{axis}[
    ybar,
    bar width=5pt,
    width=\linewidth,
    height=5.2cm,
    xmin=0.5,
    xmax=8.5,
    ymin=-1.0,
    ymax=1.0,
    axis lines=left,
    xtick={1,2,3,4,5,6,7,8},
    ytick={-1,-0.5,0,0.5,1},
    xlabel={Lag},
    ylabel={Corr.},
    enlarge x limits=0.08,
    enlarge y limits=0.12,
    clip=false
]
\addplot[draw=blue!70, fill=blue!40] coordinates {#2};
\addplot[densely dashed, gray] coordinates {(0.5,0.2) (8.5,0.2)};
\addplot[densely dashed, gray] coordinates {(0.5,-0.2) (8.5,-0.2)};
\addplot[black, thick] coordinates {(0.5,0) (8.5,0)};
\end{axis}
\end{tikzpicture}
\caption{ACF}
\end{subfigure}
\hfill
\begin{subfigure}[t]{0.48\textwidth}
\centering
\begin{tikzpicture}
\begin{axis}[
    ybar,
    bar width=5pt,
    width=\linewidth,
    height=5.2cm,
    xmin=0.5,
    xmax=8.5,
    ymin=-1.0,
    ymax=1.0,
    axis lines=left,
    xtick={1,2,3,4,5,6,7,8},
    ytick={-1,-0.5,0,0.5,1},
    xlabel={Lag},
    ylabel={Corr.},
    enlarge x limits=0.08,
    enlarge y limits=0.12,
    clip=false
]
\addplot[draw=red!70, fill=red!40] coordinates {#3};
\addplot[densely dashed, gray] coordinates {(0.5,0.2) (8.5,0.2)};
\addplot[densely dashed, gray] coordinates {(0.5,-0.2) (8.5,-0.2)};
\addplot[black, thick] coordinates {(0.5,0) (8.5,0)};
\end{axis}
\end{tikzpicture}
\caption{PACF}
\end{subfigure}

\caption{#1}
\end{figure}

\noindent #4

\medskip

#5

\medskip
}

\begin{document}

\begin{center}
{\LARGE \textbf{ACF and PACF Schematic Notes}}\\[4pt]
{\large Side-by-side illustrations for common time-series models}
\end{center}

\section*{How to read the plots}
These figures are schematic, not data-driven. They show the usual textbook patterns used for model identification.

\begin{itemize}
  \item For a stationary series, the ACF and PACF usually die down reasonably fast.
  \item For an AR model, the PACF cuts off after lag $p$ while the ACF tails off.
  \item For an MA model, the ACF cuts off after lag $q$ while the PACF tails off.
  \item For an ARMA model, both ACF and PACF tail off.
  \item For an ARIMA model, the original series is typically nonstationary, but the differenced series behaves like an ARMA model.
\end{itemize}

\section*{1. White noise}
White noise has no meaningful autocorrelation. Both plots should look like random small spikes around zero.

\schematicpair
{White noise: both ACF and PACF stay near zero}
{(1,0.05) (2,-0.04) (3,0.03) (4,-0.02) (5,0.04) (6,-0.03) (7,0.02) (8,-0.01)}
{(1,0.02) (2,-0.03) (3,0.01) (4,0.04) (5,-0.02) (6,0.03) (7,-0.01) (8,0.02)}
{Interpretation: no clear cutoff and no trend; almost all spikes stay inside the bounds.}
{This is the visual pattern of a white-noise series.}

\section*{2. AR models}
For an AR model, the ACF tails off and the PACF cuts off after the order of the model.

\subsection*{AR(1)}
\schematicpair
{AR(1): ACF tails off, PACF cuts off after lag 1}
{(1,0.70) (2,0.48) (3,0.33) (4,0.23) (5,0.16) (6,0.11) (7,0.08) (8,0.05)}
{(1,0.72) (2,0.04) (3,-0.03) (4,0.02) (5,-0.02) (6,0.01) (7,-0.01) (8,0.01)}
{Interpretation: the ACF decays gradually, but the PACF shows one dominant spike at lag 1 and then becomes small.}
{This is the usual pattern for AR(1).}

\subsection*{AR(2)}
\schematicpair
{AR(2): ACF tails off, PACF cuts off after lag 2}
{(1,0.60) (2,0.42) (3,0.25) (4,0.16) (5,0.10) (6,0.07) (7,0.05) (8,0.03)}
{(1,0.58) (2,-0.41) (3,0.03) (4,-0.02) (5,0.01) (6,-0.01) (7,0.00) (8,0.01)}
{Interpretation: the PACF shows significant lags 1 and 2, then becomes small; the ACF dies down gradually.}
{This is the usual pattern for AR(2).}

\section*{3. MA models}
For an MA model, the ACF cuts off after the order of the model and the PACF tails off.

\subsection*{MA(1)}
\schematicpair
{MA(1): ACF cuts off after lag 1, PACF tails off}
{(1,0.68) (2,0.02) (3,-0.01) (4,0.01) (5,-0.02) (6,0.00) (7,0.01) (8,-0.01)}
{(1,0.70) (2,0.35) (3,0.20) (4,0.12) (5,0.08) (6,0.05) (7,0.03) (8,0.02)}
{Interpretation: the ACF has one large spike at lag 1 and then becomes small; the PACF fades away slowly.}
{This is the usual pattern for MA(1).}

\subsection*{MA(2)}
\schematicpair
{MA(2): ACF cuts off after lag 2, PACF tails off}
{(1,0.62) (2,0.44) (3,0.03) (4,-0.02) (5,0.01) (6,0.00) (7,0.01) (8,-0.01)}
{(1,0.60) (2,0.34) (3,0.22) (4,0.15) (5,0.10) (6,0.06) (7,0.03) (8,0.02)}
{Interpretation: the first two ACF lags are the main ones, and the PACF decays gradually.}
{This is the usual pattern for MA(2).}

\section*{4. ARMA models with different $p,q$ values}
For ARMA models, neither the ACF nor the PACF has a clean cutoff. Both tend to tail off.

\subsection*{ARMA(1,1)}
\schematicpair
{ARMA(1,1): both ACF and PACF tail off}
{(1,0.58) (2,0.39) (3,0.26) (4,0.18) (5,0.12) (6,0.08) (7,0.05) (8,0.03)}
{(1,0.55) (2,0.31) (3,0.18) (4,0.11) (5,0.07) (6,0.04) (7,0.02) (8,0.01)}
{Interpretation: both plots decay gradually and no sharp cutoff appears.}
{This is typical of ARMA(1,1).}

\subsection*{ARMA(2,1)}
\schematicpair
{ARMA(2,1): both plots tail off with a short-lag structure}
{(1,0.66) (2,0.47) (3,0.28) (4,0.18) (5,0.12) (6,0.08) (7,0.05) (8,0.03)}
{(1,0.60) (2,-0.28) (3,0.14) (4,0.08) (5,0.05) (6,0.03) (7,0.02) (8,0.01)}
{Interpretation: short-lag dependence is visible, but both ACF and PACF still taper off rather than cutting off cleanly.}
{This is typical of ARMA(2,1).}

\subsection*{ARMA(1,2)}
\schematicpair
{ARMA(1,2): both plots tail off with a more flexible short-lag shape}
{(1,0.63) (2,0.41) (3,0.25) (4,0.16) (5,0.10) (6,0.06) (7,0.04) (8,0.02)}
{(1,0.58) (2,0.29) (3,-0.15) (4,0.09) (5,0.05) (6,0.03) (7,0.02) (8,0.01)}
{Interpretation: a mixed short-lag pattern appears, but neither plot has a hard cutoff.}
{This is typical of ARMA(1,2).}

\section*{5. ARIMA models with different $p,q,d$ values}
For ARIMA models, the original series is often nonstationary. After differencing $d$ times, the differenced series behaves like an ARMA model.

\subsection*{ARIMA(1,1,0)}
\schematicpair
{ARIMA(1,1,0): original series is nonstationary, differenced series looks AR(1)}
{(1,0.88) (2,0.79) (3,0.70) (4,0.63) (5,0.57) (6,0.51) (7,0.46) (8,0.41)}
{(1,0.72) (2,0.03) (3,-0.02) (4,0.01) (5,-0.02) (6,0.01) (7,0.00) (8,-0.01)}
{Interpretation: the original series has a slowly decaying ACF, while the differenced series shows a clear AR(1)-type PACF cutoff.}
{This is a common ARIMA(1,1,0) pattern.}

\subsection*{ARIMA(0,1,1)}
\schematicpair
{ARIMA(0,1,1): original series is nonstationary, differenced series looks MA(1)}
{(1,0.86) (2,0.78) (3,0.69) (4,0.62) (5,0.56) (6,0.50) (7,0.45) (8,0.40)}
{(1,0.66) (2,0.02) (3,-0.01) (4,0.01) (5,-0.01) (6,0.00) (7,0.01) (8,-0.01)}
{Interpretation: the original series is nonstationary, but after first differencing the ACF shows one dominant lag and then fades.}
{This is a common ARIMA(0,1,1) pattern.}

\subsection*{ARIMA(1,1,1)}
\schematicpair
{ARIMA(1,1,1): original series is nonstationary, differenced series looks ARMA(1,1)}
{(1,0.87) (2,0.78) (3,0.69) (4,0.62) (5,0.56) (6,0.50) (7,0.45) (8,0.40)}
{(1,0.58) (2,0.30) (3,0.17) (4,0.10) (5,0.06) (6,0.04) (7,0.02) (8,0.01)}
{Interpretation: the differenced series has short-lag dependence in both plots, but no clear cutoff in either one.}
{This is a common ARIMA(1,1,1) pattern.}

\subsection*{ARIMA(2,2,1)}
\schematicpair
{ARIMA(2,2,1): original series needs second differencing, then looks like ARMA(2,1)}
{(1,0.93) (2,0.87) (3,0.80) (4,0.74) (5,0.69) (6,0.64) (7,0.59) (8,0.54)}
{(1,0.61) (2,-0.26) (3,0.13) (4,0.08) (5,0.05) (6,0.03) (7,0.02) (8,0.01)}
{Interpretation: the raw series is strongly nonstationary, but after two differences the dependence pattern is consistent with an ARMA model.}
{This is a common ARIMA(2,2,1) pattern.}

\section*{Short summary}
\begin{center}
\begin{tabular}{@{}lll@{}}
\toprule
Model class & ACF pattern & PACF pattern \\
\midrule
White noise & no structure & no structure \\
AR($p$) & tails off & cuts off after $p$ \\
MA($q$) & cuts off after $q$ & tails off \\
ARMA($p,q$) & tails off & tails off \\
ARIMA($p,d,q$) & differenced series follows ARMA($p,q$) & differenced series follows ARMA($p,q$) \\
\bottomrule
\end{tabular}
\end{center}

\end{document}
